%%%%%%%%%%%%%%%%%%%%%%%%%%%%%%%%%%%%%%%%%%%%%%%%%%%%%%%%%%%%%%%%%%%%%%%%%%%%%%%%%%
\begin{frame}[fragile]\frametitle{}
\begin{center}
{\Large Prompt Engineering for Everyday Work}

\medskip
{\large Office Professionals \& Teachers}
\end{center}
\end{frame}


%%%%%%%%%%%%%%%%%%%%%%%%%%%%%%%%%%%%%%%%%%%%%%%%%%%%%%%%%%%
\begin{frame}[fragile]\frametitle{Why This Matters for You}

\begin{columns}
    \begin{column}[T]{0.55\linewidth}
    You do not need to be a programmer to benefit from Prompt Engineering.

    \medskip
    \textbf{These are real tasks you do every week:}
    \begin{itemize}
        \item Writing emails you keep putting off
        \item Preparing reports and presentations
        \item Summarising long documents or meeting notes
        \item Replying to difficult or sensitive messages
        \item Planning lessons or training sessions
        \item Giving feedback on student or colleague work
        \item Organising your own thoughts before a meeting
    \end{itemize}
    \end{column}
    \begin{column}[T]{0.45\linewidth}
    \begin{center}
    \renewcommand{\arraystretch}{1.4}
    \begin{tabular}{|l|c|}
    \hline
    \textbf{Task} & \textbf{Time Saved} \\
    \hline
    Drafting email & 80\% \\
    Meeting summary & 70\% \\
    Report first draft & 65\% \\
    Lesson plan & 60\% \\
    Feedback writing & 75\% \\
    \hline
    \end{tabular}
    \end{center}
    \medskip
    {\tiny Estimates based on typical knowledge worker time studies.}
	
	    \textbf{A well-crafted prompt can cut each of these tasks from 30 minutes to 3 minutes.}
    \end{column}
\end{columns}

\end{frame}


%%%%%%%%%%%%%%%%%%%%%%%%%%%%%%%%%%%%%%%%%%%%%%%%%%%%%%%%%%%
\begin{frame}[fragile]\frametitle{Writing Emails You Keep Avoiding}

    \begin{itemize}
        \item \textbf{Situation:} You need to email a client about a delayed project. You have been avoiding it for two days.
        \item \textbf{Key technique:} Bullet-point your requirements. The AI respects every constraint simultaneously.
    \end{itemize}

\begin{lstlisting}[language=HTML]
Prompt:
Write a professional email to a client informing them that their software project will be delayed by 2 weeks  due to unexpected technical issues discovered during testing. 
- Tone: apologetic but confident
- Acknowledge the inconvenience
- Briefly explain the reason (without too much jargon)
- Reassure them the delay ensures a better final product
- Offer a call to discuss if needed
- Keep it under 150 words

Output:
Subject: Update on Your Project Timeline

Dear [Client Name],
I wanted to reach out personally regarding your project. During our final testing phase, we identified some  technical issues that need to be resolved properly...
\end{lstlisting}


\end{frame}


%%%%%%%%%%%%%%%%%%%%%%%%%%%%%%%%%%%%%%%%%%%%%%%%%%%%%%%%%%%
\begin{frame}[fragile]\frametitle{Replying to a Difficult or Angry Message}

    \begin{itemize}
        \item \textbf{Situation:} You received a sharp email from a colleague or parent/guardian. You are irritated and do not want to write back rashly.
        \item \textbf{Why this works:} You offload the emotional labour of finding the right words to the AI. You review, adjust slightly, and send, rather than stewing for hours.
        \item 	\textbf{Pro tip:} Always read the output carefully and personalise it. The AI gives you the structure; your judgment adds the authenticity.
    \end{itemize}
		
\begin{lstlisting}[language=HTML]
Prompt:
I received this email from an angry parent:
"""
[Paste the email here]
"""
Write a calm, professional reply that:
- Acknowledges their concern without being defensive
- Does not admit fault, but shows empathy
- Proposes a 20-minute call this week to resolve it
- Ends on a positive, collaborative note
- Tone: warm but firm
\end{lstlisting}




\end{frame}


%%%%%%%%%%%%%%%%%%%%%%%%%%%%%%%%%%%%%%%%%%%%%%%%%%%%%%%%%%%
\begin{frame}[fragile]\frametitle{Summarising Long Documents and Reports}

    \begin{itemize}
        \item \textbf{Situation:} You have a 40-page policy document, a lengthy meeting transcript, or a research paper to read before a meeting in one hour.
        \item \textbf{Output format tip:} Asking for a \emph{numbered list of specific things} is far more useful than asking to ``summarise this'', you get exactly the information you need, not a generic paragraph.
        \item \textbf{Works for:} Annual reports, HR policies, legal notices, research papers, long email threads.
    \end{itemize}
		
\begin{lstlisting}[language=HTML]
Prompt:
Below is a document. Read it carefully and give me:
1. A 5-sentence executive summary
2. The 3 most important decisions or changes being proposed
3. Any deadlines or dates mentioned
4. One question I should ask in the meeting about this

"""
[Paste the document text here]
"""
\end{lstlisting}



\end{frame}


%%%%%%%%%%%%%%%%%%%%%%%%%%%%%%%%%%%%%%%%%%%%%%%%%%%%%%%%%%%
\begin{frame}[fragile]\frametitle{Converting Meeting Notes into Action Items}

    \begin{itemize}
        \item \textbf{Situation:} You took rough notes during a 1-hour meeting. Now you need to share a clean summary with the team.
        \item \textbf{Insight:} You can paste rough, messy notes, the AI does not judge grammar or spelling. It extracts structure from chaos.
    \end{itemize}

\begin{lstlisting}[language=HTML]
Prompt:
These are my rough meeting notes from today:
"""
[Paste your rough notes, typos and all]
"""

Please:
1. Write a clean 3-paragraph meeting summary
2. List all action items in a table with columns: 
   Action | Owner | Deadline
3. List any unresolved questions that need follow-up
4. Write a short 4-line email I can send to attendees 
   with the summary attached

Output format: Use clear headings for each section.
\end{lstlisting}


\end{frame}


%%%%%%%%%%%%%%%%%%%%%%%%%%%%%%%%%%%%%%%%%%%%%%%%%%%%%%%%%%%
\begin{frame}[fragile]\frametitle{Writing Performance Reviews and Feedback}

    \begin{itemize}
        \item \textbf{Situation:} You need to write a performance review for a team member or a progress report on a student. You want it to be fair, constructive, and specific.
        \item \textbf{Key technique:} Say explicitly what to \emph{avoid}. ``Avoid vague phrases'' produces dramatically more useful feedback than a generic request.
    \end{itemize}

\begin{lstlisting}[language=HTML]
Prompt:
Write a constructive performance review paragraph for 
a team member with the following profile:
- Role: Junior Data Analyst
- Strengths: very detail-oriented, meets deadlines, 
  proactive in asking questions
- Areas to improve: needs to communicate findings more 
  clearly to non-technical stakeholders; sometimes 
  over-explains technical details in meetings
- Overall: solid performer, on track for promotion 
  in 12 months

Tone: encouraging, specific, professional.
Length: 150-180 words.
Avoid vague phrases like "good team player."
\end{lstlisting}



\end{frame}


%%%%%%%%%%%%%%%%%%%%%%%%%%%%%%%%%%%%%%%%%%%%%%%%%%%%%%%%%%%
\begin{frame}[fragile]\frametitle{Creating a Lesson Plan (For Teachers)}

    \begin{itemize}
        \item \textbf{Situation:} You need to plan a 45-minute lesson on a topic for a specific age group, today.
        \item \textbf{Result:} A complete, structured, ready-to-use lesson plan in under 30 seconds. Adapt the output to your classroom's specific needs.
    \end{itemize}
		
\begin{lstlisting}[language=HTML]
Prompt:
Create a detailed 45-minute lesson plan for:
- Topic: Introduction to Fractions
- Grade: 5th standard (age 10-11)
- Prior knowledge: students know division basics
- Learning objective: understand what a fraction means 
  and compare simple fractions (1/2, 1/4, 3/4)

Include:
- 5-minute warm-up activity (no materials needed)
- Main concept explanation with a real-life analogy
- A hands-on activity using items in a classroom
- 5 practice problems (easy to hard)
- A 2-minute closing question to check understanding
- One common misconception to watch out for
\end{lstlisting}



\end{frame}


%%%%%%%%%%%%%%%%%%%%%%%%%%%%%%%%%%%%%%%%%%%%%%%%%%%%%%%%%%%
\begin{frame}[fragile]\frametitle{Generating Quiz and Exam Questions}

    \begin{itemize}
        \item \textbf{Situation:} You need to create 20 questions for an assessment. Writing good questions is time-consuming, especially writing good \emph{wrong} answers for multiple-choice.
        \item \textbf{Why ``plausible wrong options'' matters:} Vague wrong answers make questions trivially easy. Asking the AI to use common misconceptions creates genuinely diagnostic questions.
    \end{itemize}
		
\begin{lstlisting}[language=HTML]
Prompt:
Create 10 multiple-choice questions on the topic: 
"The Water Cycle" for Grade 7 students.

For each question:
- Write one clear question
- Provide 4 options (A, B, C, D)
- Make sure the 3 wrong options are plausible 
  (common misconceptions, not obviously silly)
- Mark the correct answer
- Add a one-line explanation of why it is correct

Difficulty: mix of easy (4), medium (4), hard (2).
\end{lstlisting}



\end{frame}


%%%%%%%%%%%%%%%%%%%%%%%%%%%%%%%%%%%%%%%%%%%%%%%%%%%%%%%%%%%
\begin{frame}[fragile]\frametitle{Explaining Difficult Concepts at Different Levels}

    \begin{itemize}
        \item \textbf{Situation:} You need to explain the same concept to different audiences, a student, a parent, and the school principal.
        \item \textbf{Why this is powerful:} You write the concept once; the AI adapts the vocabulary, analogies, and tone for each audience automatically.
        \item \textbf{Office use:} Same technique for explaining a technical issue to your manager vs. to an IT colleague vs. to a client.
    \end{itemize}

\begin{lstlisting}[language=HTML]
Prompt:
Explain the concept of "compound interest" three times:
1. To a 12-year-old student using a pocket money analogy
2. To a parent at a school financial literacy evening, 
   practical, relatable, no jargon
3. To a school board member in one precise paragraph 
   for an annual report

Keep each explanation self-contained and appropriate 
for the audience.
\end{lstlisting}


\end{frame}


%%%%%%%%%%%%%%%%%%%%%%%%%%%%%%%%%%%%%%%%%%%%%%%%%%%%%%%%%%%
\begin{frame}[fragile]\frametitle{Drafting Proposals and Business Cases}

    \begin{itemize}
        \item \textbf{Situation:} You need to propose a new initiative, a budget request, a process change, or a new school programme, and need to make it persuasive.
        \item \textbf{Key technique:} Pre-answering objections in the prompt forces the AI to make a stronger, more complete case.
    \end{itemize}
		
\begin{lstlisting}[language=HTML]
Prompt:
Write a one-page business case for introducing a 
"No-Meeting Wednesday" policy in our 80-person 
software company.

Structure:
- Problem: meeting overload reducing deep work time
- Proposed Solution: one meeting-free day per week
- Expected Benefits: 3 specific, measurable outcomes
- Potential Objections: 2 objections + how to address them
- Simple 4-week pilot plan
- Recommended next step (one sentence call to action)

Tone: professional, evidence-based, concise.
Audience: Senior leadership team.
\end{lstlisting}



\end{frame}


%%%%%%%%%%%%%%%%%%%%%%%%%%%%%%%%%%%%%%%%%%%%%%%%%%%%%%%%%%%
\begin{frame}[fragile]\frametitle{Preparing for a Difficult Conversation}

    \begin{itemize}
        \item \textbf{Situation:} You have to give critical feedback to a colleague, negotiate with a vendor, or handle a parent complaint. You want to prepare what to say.
        \item \textbf{Insight:} The AI becomes your \emph{personal communication coach}, available any time, at zero cost.
    \end{itemize}

		
\begin{lstlisting}[language=HTML]
Prompt:
I need to have a conversation with a colleague who 
consistently misses deadlines, affecting the whole team.
This is the third time in two months.

Help me prepare by giving me:
1. An opening sentence that is direct but not aggressive
2. The 3 key points I must cover
3. How to phrase the core issue without using 
   accusatory language ("you always..." style)
4. How to respond if they become defensive
5. A closing statement that ends on a constructive note

My goal: agree on a clear plan going forward, not 
just vent my frustration.
\end{lstlisting}


\end{frame}


%%%%%%%%%%%%%%%%%%%%%%%%%%%%%%%%%%%%%%%%%%%%%%%%%%%%%%%%%%%
\begin{frame}[fragile]\frametitle{Creating Training and Onboarding Materials}

    \begin{itemize}
        \item \textbf{Situation:} A new employee or student is joining. You need to create a guide, checklist, or FAQ quickly.
        \item \textbf{Time saved:} What normally takes a senior staff member an afternoon of writing now takes 2 minutes to draft, 10 minutes to review and customise.
    \end{itemize}
		
\begin{lstlisting}[language=HTML]
Prompt:
Create a "First Week Onboarding Checklist" for a new 
primary school teacher joining mid-year.

Include:
- Day 1: Essential admin tasks (HR forms, ID card, etc.)
- Day 2-3: Classroom setup and resource access
- Day 4-5: Meeting key people (checklist format)
- One-page FAQ answering the 8 most common questions 
  a new teacher asks (e.g., how to raise a discipline 
  issue, how to request supplies, marking policy)
- A "Who to Call" quick reference card

Format everything so it can be printed on 2 A4 pages.
\end{lstlisting}



\end{frame}


%%%%%%%%%%%%%%%%%%%%%%%%%%%%%%%%%%%%%%%%%%%%%%%%%%%%%%%%%%%
\begin{frame}[fragile]\frametitle{Analysing Survey or Feedback Data}

    \begin{itemize}
        \item \textbf{Situation:} You collected feedback forms, from students, parents, or clients, and need to make sense of 50 free-text responses.
        \item \textbf{What would take hours} of reading, categorising, and writing is done in seconds. You validate and add context, the AI does the heavy lifting.
    \end{itemize}
		
\begin{lstlisting}[language=HTML]
Prompt:
Below are 50 responses from our annual parent 
satisfaction survey. Read all of them carefully.

"""
[Paste all responses here]
"""

Then provide:
1. Top 5 themes that appear most frequently 
   (with example quotes for each)
2. Top 3 specific complaints (most actionable)
3. Top 3 specific compliments (worth celebrating)
4. One surprising or unexpected insight
5. A short paragraph I can include in my report 
   summarising overall sentiment (positive/neutral/
   negative split as a rough percentage)
\end{lstlisting}



\end{frame}


%%%%%%%%%%%%%%%%%%%%%%%%%%%%%%%%%%%%%%%%%%%%%%%%%%%%%%%%%%%
\begin{frame}[fragile]\frametitle{Writing Social Media and Newsletter Content}

    \begin{itemize}
        \item \textbf{Situation:} Your school or office needs regular communication, a monthly newsletter, a LinkedIn post, or a WhatsApp group update. Nobody has time to write it from scratch.
        \item \textbf{Technique:} Asking for multiple versions lets you pick the best one or mix elements from different versions. Always faster than writing even one version from scratch.
    \end{itemize}
		
\begin{lstlisting}[language=HTML]
Prompt:
Write 3 versions of a LinkedIn post announcing that 
our school won the District Science Fair for the 
second consecutive year.

Version 1: Formal and proud (for school's official page)
Version 2: Warm and celebratory (for the principal's 
           personal profile)
Version 3: Energetic and fun, targeting students 
           and young alumni

Each version: 3-4 sentences, include a call to action 
(e.g., congratulate the students in comments).
Add 5 relevant hashtags for each version.
\end{lstlisting}



\end{frame}


%%%%%%%%%%%%%%%%%%%%%%%%%%%%%%%%%%%%%%%%%%%%%%%%%%%%%%%%%%%
\begin{frame}[fragile]\frametitle{Building a Personal ``Prompt Library''}

\textbf{The highest-leverage habit you can build:}

\medskip
Once you find a prompt that works well for a recurring task, \textbf{save it as a reusable template}.

\medskip
Replace the \texttt{[PLACEHOLDERS]} each time. \textbf{Your prompt library becomes your personal AI productivity toolkit.}

\begin{lstlisting}[language=HTML]
Example saved prompt templates:

[EMAIL_DELAY]
Write a professional email informing [CLIENT] that 
[PROJECT] is delayed by [DURATION] due to [REASON]. 
Tone: apologetic but confident. Under 150 words.

[LESSON_PLAN]  
Create a 45-minute lesson plan for [TOPIC], Grade [X], 
with prior knowledge of [Y]. Include warm-up, activity, 
5 practice problems, and a closing check question.

[MEETING_SUMMARY]
From these rough notes [NOTES], write: a 3-para summary, 
an action table (Action | Owner | Deadline), and a 
4-line email to attendees.
\end{lstlisting}



\end{frame}


%%%%%%%%%%%%%%%%%%%%%%%%%%%%%%%%%%%%%%%%%%%%%%%%%%%%%%%%%%%
\begin{frame}[fragile]\frametitle{The Golden Rules for Everyday Prompting}

\begin{columns}
    \begin{column}[T]{0.5\linewidth}
    \textbf{Always include:}
    \begin{itemize}
        \item \textbf{Context}, who you are, what the situation is
        \item \textbf{Audience}, who will read the output
        \item \textbf{Tone}, formal / friendly / assertive
        \item \textbf{Format}, bullet list / table / email / paragraph
        \item \textbf{Length}, word count or ``under X sentences''
        \item \textbf{Constraints}, what to avoid
    \end{itemize}
    \end{column}
    \begin{column}[T]{0.5\linewidth}
    \textbf{Always remember:}
    \begin{itemize}
        \item \textbf{First draft, not final output}, always review and personalise
        \item \textbf{Iterate}, if the first output is not quite right, refine the prompt and try again
        \item \textbf{Paste real content}, the AI works best when you give it the actual text, email, or notes to work with
        \item \textbf{Never paste passwords, personal IDs, or confidential data} into public AI tools
        \item \textbf{Save what works}, build your prompt library
    \end{itemize}
    \end{column}
\end{columns}

\medskip
\begin{center}
\textbf{The best prompt is specific, contextual, and tells the AI exactly what ``done'' looks like.}
\end{center}

\end{frame}
